\documentclass{situationist-preprint}

\title{The Non-Identity Problem}
\msid{ETH-08-b}
\msfamily{ETH-08}
\msdiscipline{Normative and Applied Ethics}
\mstype{Companion manuscript}
\mskeywords{non-identity problem; procreation; fourteen-year-old girl; impersonal obligation; wronging}

\begin{document}
\maketitle

\begin{abstract}
ETH-08-a argues that obligations to future generations are impersonal: identity-affecting choices that make the future worse wrong no particular person. This manuscript applies the argument to individual procreative choices, where the non-identity problem first appeared. We examine the fourteen-year-old girl, the woman advised to delay conception, and the case of a child conceived for slavery. In the first two, the child is not wronged, and whether the parent acts wrongly depends on how the impersonal reason weighs against the cost of waiting. In the third, the child is wronged, but by acts that come after his conception. The cases confirm that procreative obligations, where there are any, are owed to no one in particular.
\end{abstract}

\section{The framework}

We rely on two principles of ETH-08-a. By the comparative condition on harming (Principle~P1), a choice harms a person only if it makes her worse off than she would otherwise have been. By the claim condition on wronging (Principle~P2), a choice wrongs a person only if it harms her or infringes a claim she holds, and it infringes such a claim only if some alternative open to the agent is one in which she exists and the claim is honored. Lemmas~3.2 and~3.3 of ETH-08-a apply these principles to policies; the same reasoning applies to a single act of conception, which by Assumption~2.4 of ETH-08-a fixes which person will exist.

The positive account is Proposition~4.1 of ETH-08-a: given the impersonal axiology of ETH-09, an identity-affecting choice is wrong, when it is, because it brings about a worse outcome when a better one was available at moderate cost. Impersonal obligations are not obligations to maximize. Whether a reason to bring about the better outcome becomes an obligation depends on the cost to the agent.

\section{Three cases}

\paragraph{The fourteen-year-old girl.} \citet[ch.~16]{Parfit1984} describes a girl who chooses to have a child at fourteen. Because she is so young, she gives her child a bad start in life; had she waited, she would have had a different child, to whom she would have given a better start. The child she has has a life worth living.

The child is not harmed: there is no alternative in which he exists and is better off, so Principle~P1 of ETH-08-a is not satisfied. Nor is he wronged. Any claim he could hold against his mother, for instance a claim to a better start, is a claim that no alternative open to her honors while he exists, since the only alternative in which he has a better start is one in which he is not conceived at all. By Principle~P2 of ETH-08-a, he has no claim that her choice infringes.

What remains is the impersonal reason. The outcome in which she waits contains a life that goes better, with the same number of people, so by the same-number principle that ETH-08-a takes from Parfit it is a better outcome. Whether she is obligated to wait depends on the cost to her. If the cost is small, she ought to wait, and in conceiving at fourteen she acts wrongly, though she wrongs no one. If the cost is considerable, for instance because waiting would frustrate a reasonable life plan, the impersonal reason may be outweighed, and then she does not act wrongly. Procreative autonomy enters at this point, not as a right that trumps all reasons, but as one of the goods of the girl's own life that the weighing must count.

This verdict explains a feature of common opinion that person-affecting accounts struggle with. Most people think the girl should wait, but few think that her child, once born, has a grievance against her. On our account both judgements are correct. The reason to wait is real, and it is not a reason grounded in the child's claims.

\paragraph{Delayed conception.} A woman learns that, because of a temporary condition, a child she conceives now will have a significant impairment, while a child conceived three months later will not. Waiting costs her little. The case has the same structure as the girl's, with the cost to the agent nearly zero. The impersonal reason is undefeated, she ought to wait, and if she does not she acts wrongly. Her child, who exists only because she did not wait, is not wronged. The verdict agrees with the No-Difference View that ETH-08-a defends: her conduct is as wrong as that of a pregnant woman who could prevent the same impairment in her existing fetus by taking a costless medicine, and wrong for the same impersonal reason, since the two outcomes differ in the same way in the quality of the lives they contain.

Some writers resist this verdict by appealing to the wrong of creating a person in a condition that it is bad to be in, whether or not a better alternative existed for that person \citep{Shiffrin1999}. Such an account must hold that the impaired child is wronged by being brought into existence with a life worth living. We accept that bringing someone into existence with an impairment can impose burdens that call for compensation. But a burden imposed in the only course in which its bearer exists, and in a life she has reason to be glad of, is not one that grounds a claim against the agent under Principle~P2.

\paragraph{The slave child.} \citet{Kavka1982} describes a couple who, in return for payment, conceive a child to be raised as a slave. The child's life, we may suppose, is worth living, and he would not have existed had the couple refused the bargain. The case has been taken to show that the non-identity problem proves too much: surely the child is wronged.

He is, but not by his conception. The couple's later acts, and the acts of his owner, are not identity-affecting: once the child exists, there are alternatives in which he exists and is not enslaved, and in which his claims to liberty are honored. By Principle~P2 of ETH-08-a, the enslavement infringes his claims, and since it is impermissible it violates them. The bargain is wrong in two ways at once. It is wrong because it commits the couple to acts that will wrong the child once he exists, and it is wrong impersonally, as a choice of a worse outcome. What it is not, considered as an act of conception, is a wrong to the child. This is the pattern ETH-08-a found in Buried Glass: obligations concerning people who will exist whatever we do, once the identity-fixing choice has been made, are directed.

\section{Objection and reply}

\begin{objection}
An act can wrong a person by infringing his rights even when it does not make him worse off \citep{Woodward1986}. The girl's child has a right to an adequate start in life, and by conceiving at fourteen she foreseeably infringes it.
\end{objection}

\begin{reply}
We accept the first sentence, which Principle~P2 of ETH-08-a also accepts: harm is not required for wronging. The question is whether the right in the second sentence is one the girl's choice infringes. If it is a right to care that she can give once the child exists, then there are alternatives in which he exists and it is honored, and her later failures to give it, if any, wrong him, as in the slave case. If it is a right to have been conceived under better conditions, no option open to her honors it while he exists, since conceiving later would have produced someone else, and by Principle~P2 of ETH-08-a it is not infringed. The reason against her choice that the objection registers is real, but it is the impersonal reason that a better outcome was available. ETH-08-a reaches the same verdict about rights-based replies in general in its Section~5.
\end{reply}

\section{Conclusion}

In individual procreative choices, as in policies, the child brought into existence by an identity-affecting choice is not wronged by that choice. The reason against conceiving under bad conditions is impersonal, and whether it becomes an obligation depends on what the alternative costs the parent. Where a child is wronged, as in the slave case, the wrong is done by later acts that are not identity-affecting. The cases support the thesis of ETH-08-a that our obligations concerning who will exist and how well they will live are owed to no one in particular.

\begin{thebibliography}{9}\small
\bibitem[Kavka(1982)]{Kavka1982} Kavka, G.~S. (1982). The paradox of future individuals. \emph{Philosophy \& Public Affairs}, 11(2), 93--112.
\bibitem[Parfit(1984)]{Parfit1984} Parfit, D. (1984). \emph{Reasons and Persons}. Clarendon Press.
\bibitem[Shiffrin(1999)]{Shiffrin1999} Shiffrin, S.~V. (1999). Wrongful life, procreative responsibility, and the significance of harm. \emph{Legal Theory}, 5(2), 117--148.
\bibitem[Woodward(1986)]{Woodward1986} Woodward, J. (1986). The non-identity problem. \emph{Ethics}, 96(4), 804--831.
\end{thebibliography}

\end{document}
